\documentclass[12pt,a4paper]{article}

% --- Metadata (per course) ---
\def\courseshortheader{QUÁ TRÌNH NGẪU NHIÊN}
\def\coverbookline{NHÓM 5 --- XÁC SUẤT VÀ THỐNG KÊ}
\def\covermaintitle{Quá trình ngẫu nhiên}
\def\coversubtitle{Giáo án lý thuyết --- Markov, Poisson, martingale, Brown}

% Shared preamble for nhom_5 (requires \courseshortheader before \input)
\usepackage[utf8]{inputenc}
\usepackage[vietnamese]{babel}
\usepackage{amsmath,amssymb,amsthm}
\usepackage{geometry}
\usepackage{enumitem}
\usepackage[most]{tcolorbox}
\usepackage{hyperref}
\usepackage{fancyhdr}
\usepackage{graphicx}
\usepackage{booktabs}
\usepackage{tikz}
\usetikzlibrary{calc,arrows.meta,decorations.pathreplacing,positioning}
\usepackage{eso-pic}
\usepackage{xcolor}

\geometry{a4paper, margin=2cm, bottom=2.5cm}
\hypersetup{colorlinks=true, linkcolor=blue!70!black, urlcolor=blue}
\setlist[itemize]{topsep=4pt, itemsep=2pt}
\setlist[enumerate]{topsep=4pt, itemsep=2pt}

\AddToShipoutPictureFG{%
  \AtPageCenter{%
    \begin{tikzpicture}[remember picture, overlay]
      \node[rotate=45, scale=1, opacity=0.06]{%
        \fontsize{60}{72}\selectfont\bfseries\sffamily Đặng Tấn Quí%
      };
    \end{tikzpicture}%
  }%
}

\pagestyle{fancy}
\fancyhf{}
\fancyhead[L]{\small\textit{\courseshortheader}}
\fancyhead[R]{\small\textit{Đặng Tấn Quí}}
\fancyfoot[C]{\thepage}
\fancyfoot[L]{\footnotesize\textcopyright\ Đặng Tấn Quí}
\fancyfoot[R]{\footnotesize SĐT: 0377\,122\,210}
\renewcommand{\headrulewidth}{0.4pt}
\renewcommand{\footrulewidth}{0.4pt}

\fancypagestyle{plain}{%
  \fancyhf{}
  \fancyfoot[C]{\thepage}
  \fancyfoot[L]{\footnotesize\textcopyright\ Đặng Tấn Quí}
  \fancyfoot[R]{\footnotesize SĐT: 0377\,122\,210}
  \renewcommand{\headrulewidth}{0pt}
  \renewcommand{\footrulewidth}{0.4pt}
}

\newtcolorbox{dongco}[1][]{%
  enhanced, breakable,
  colback=teal!4!white, colframe=teal!60!black,
  fonttitle=\bfseries, title={Động cơ học -- Tại sao cần khái niệm này?}, #1%
}
\newtcolorbox{ynghia}[1][]{%
  enhanced, breakable,
  colback=purple!3!white, colframe=purple!50!black,
  fonttitle=\bfseries, title={Ý nghĩa \& Trực giác}, #1%
}
\newtcolorbox{ungdung}[1][]{%
  enhanced, breakable,
  colback=olive!5!white, colframe=olive!60!black,
  fonttitle=\bfseries, title={Ứng dụng thực tế}, #1%
}
\newtcolorbox{dinhnghia}[1][]{%
  enhanced, breakable,
  colback=blue!3!white, colframe=blue!60!black,
  fonttitle=\bfseries, title={Định nghĩa}, #1%
}
\newtcolorbox{dinhly}[1][]{%
  enhanced, breakable,
  colback=green!3!white, colframe=green!50!black,
  fonttitle=\bfseries, title={Định lý}, #1%
}
\newtcolorbox{tinhchat}[1][]{%
  enhanced, breakable,
  colback=yellow!5!white, colframe=orange!50!black,
  fonttitle=\bfseries, title={Tính chất}, #1%
}
\newtcolorbox{phuongphap}[1][]{%
  enhanced, breakable,
  colback=gray!5!white, colframe=gray!50!black,
  fonttitle=\bfseries, title={Phương pháp}, #1%
}
\newtcolorbox{luuy}[1][]{%
  enhanced, breakable,
  colback=red!3!white, colframe=red!50!black,
  fonttitle=\bfseries, title={Lưu ý}, #1%
}
\newtcolorbox{congthuc}[1][]{%
  enhanced, breakable,
  colback=cyan!3!white, colframe=cyan!50!black,
  fonttitle=\bfseries, title={Công thức}, #1%
}
\newtcolorbox{vidu}[1][]{%
  enhanced, breakable,
  colback=violet!3!white, colframe=violet!50!black,
  fonttitle=\bfseries, title={Ví dụ}, #1%
}
\newtcolorbox{phanvidu}[1][]{%
  enhanced, breakable,
  colback=red!2!white, colframe=red!40!black,
  fonttitle=\bfseries, title={Phản ví dụ}, #1%
}
\newtcolorbox{tongket}[1][]{%
  enhanced, breakable,
  colback=blue!4!white, colframe=blue!70!black,
  fonttitle=\bfseries, title={Tổng kết chương}, #1%
}


\begin{document}

\thispagestyle{empty}
\begin{tikzpicture}[remember picture, overlay]
  \fill[blue!80!black] (current page.south west) rectangle (current page.north east);
  \fill[blue!60!cyan, opacity=0.8]
    (current page.south west) -- (current page.north west) --
    ([xshift=-3cm]current page.north east) -- (current page.south east) -- cycle;
  \foreach \r/\o in {6/0.05, 4.5/0.08, 3/0.12} {
    \fill[white, opacity=\o] ([xshift=2cm, yshift=-2cm]current page.north east) circle (\r cm);
  }
  \draw[white, line width=2pt] ([yshift=-5cm]current page.north west) ++(2cm,0) -- ++(17cm,0);
  \draw[white, line width=2pt] ([yshift=-16cm]current page.north west) ++(2cm,0) -- ++(17cm,0);
  \node[anchor=west, white, font=\fontsize{28}{34}\bfseries\selectfont]
    at ([xshift=3cm, yshift=-7.5cm]current page.north west) {\coverbookline};
  \node[anchor=west, white, font=\fontsize{20}{26}\selectfont]
    at ([xshift=3cm, yshift=-9cm]current page.north west) {\covermaintitle};
  \node[anchor=west, white, font=\fontsize{16}{22}\selectfont]
    at ([xshift=3cm, yshift=-10.2cm]current page.north west) {\coversubtitle};
  \node[anchor=west, white, font=\Large\bfseries]
    at ([xshift=3cm, yshift=-13cm]current page.north west) {Biên soạn:};
  \node[anchor=west, yellow!80!white, font=\fontsize{24}{30}\bfseries\selectfont]
    at ([xshift=3cm, yshift=-14.5cm]current page.north west) {ĐẶNG TẤN QUÍ};
  \node[anchor=west, white!90, font=\large]
    at ([xshift=3cm, yshift=-18cm]current page.north west) {SĐT: 0377\,122\,210};
  \node[anchor=south, white!80, font=\small]
    at ([yshift=1.5cm]current page.south) {Tài liệu có bản quyền --- Cấm sao chép khi chưa được sự đồng ý của tác giả};
  \node[anchor=south east, white, font=\Large\bfseries]
    at ([xshift=-2cm, yshift=3cm]current page.south east) {\the\year};
\end{tikzpicture}

\newpage
\tableofcontents
\newpage

\section*{Lời nói đầu}
\addcontentsline{toc}{section}{Lời nói đầu}

Tài liệu thuộc \textbf{Nhóm 5 --- Xác suất và thống kê}, biên soạn theo cùng triết lý với giáo án Đại số tuyến tính: học để \emph{hiểu} và \emph{dùng}, không chỉ để ghi nhớ công thức.

\begin{enumerate}
  \item \textbf{Động cơ trước định nghĩa.} Mỗi phần lớn mở bằng ô \textsf{\color{teal!60!black} Động cơ học} --- câu hỏi thực tế hoặc bài toán mẫu cần khái niệm đó.
  \item \textbf{Trực giác trước công thức.} Ô \textsf{\color{purple!50!black} Ý nghĩa \& Trực giác} giải thích ``công thức đang nói gì'' (xác suất = tần suất dài hạn, kỳ vọng = trung bình có trọng số, v.v.).
  \item \textbf{Ví dụ có kiểm tra.} Ô \textsf{\color{violet!50!black} Ví dụ} nên tự làm trước khi đọc lời giải; phần thống kê ứng dụng cần gắn dữ liệu số cụ thể.
  \item \textbf{Tổng kết cuối mỗi chương.} Ô \textsf{\color{blue!70!black} Tổng kết chương} liệt kê kết quả phải nhớ và liên kết sang phần sau.
  \item \textbf{Phân tầng theo môn.} X1--X3 là nền; X4--X5 nâng cao lý thuyết; X6--X8 chuyên sâu (đa biến, Bayes, quá trình ngẫu nhiên).
\end{enumerate}

\noindent\textbf{Cách đọc:} Động cơ $\to$ Định nghĩa / Định lý $\to$ Ví dụ $\to$ Ý nghĩa $\to$ Tổng kết. Với môn thống kê, luôn hỏi: \emph{mẫu là gì, tham số nào chưa biết, giả thuyết $H_0$ là gì}.

\newpage


\section{Mở đầu}

\begin{dongco}[title={Động cơ --- Mở đầu}]
Quá trình ngẫu nhiên mô tả hệ thống \emph{thay đổi theo thời gian}: hàng đợi, giá tài sản, mạng sinh học. Khác BNN tĩnh: ta quan tâm quỹ đạo $\{X_t\}_{t \ge 0}$.
\end{dongco}

\newpage

\section{Tổng quan}

\begin{dongco}[title={Động cơ --- Tổng quan}]
Định nghĩa quá trình, phân phối hữu hạn chiều, tính Markov và độc lập gia tăng.
\end{dongco}

\begin{dinhnghia}
  \textbf{Quá trình ngẫu nhiên} (stochastic process): họ BNN $\{X_t\}_{t \in T}$ trên $(\Omega, \mathcal{F}, P)$.
  \begin{itemize}
    \item $T = \{0, 1, 2, \ldots\}$: thời gian rời rạc.
    \item $T = [0, \infty)$: thời gian liên tục.
    \item \textbf{Quỹ đạo} (trajectory): $t \mapsto X_t(\omega)$ với $\omega$ cố định.
    \item \textbf{Phân phối hữu hạn chiều:} $(X_{t_1}, \ldots, X_{t_n})$ xác định phân phối quá trình (Kolmogorov).
  \end{itemize}
\end{dinhnghia}

\begin{dinhnghia}[title={Phân loại}]
  \begin{itemize}
    \item \textbf{Dừng} (stationary): phân phối đồng thời bất biến với dịch thời gian.
    \item \textbf{Dừng yếu} (wide-sense): $E[X_t] = \mu$ const, $\text{Cov}(X_s, X_t) = R(|t-s|)$.
    \item \textbf{Độc lập gia số}: $X_{t_1} - X_{t_0}$, $X_{t_2} - X_{t_1}$, \ldots\ độc lập.
  \end{itemize}
\end{dinhnghia}

%% ================================================================
%%  CHƯƠNG 2: XÍCH MARKOV
%% ================================================================

\begin{tongket}[title={Tổng kết: Tổng quan}]
\begin{enumerate}
  \item Nắm lại các \textbf{định nghĩa} và \textbf{công thức} cốt lõi trong mục ``Tổng quan''.
  \item Tự làm lại ít nhất một ví dụ (nếu có) hoặc áp dụng công thức vào một tình huống số.
  \item Ghi chú điều kiện áp dụng (giả thiết độc lập, phân phối chuẩn, $n$ đủ lớn, v.v.).
\end{enumerate}
\end{tongket}

\section{Xích Markov}

\begin{dongco}[title={Động cơ --- Xích Markov}]
Tương lai chỉ phụ thuộc hiện tại --- mô hình hóa nhiều hệ rời rạc.
\end{dongco}

\begin{dinhnghia}
  $(X_n)_{n \ge 0}$ trên không gian trạng thái $S$ (đếm được) là \textbf{xích Markov} nếu:
  \[
    P(X_{n+1} = j \mid X_n = i, X_{n-1}, \ldots, X_0) = P(X_{n+1} = j \mid X_n = i) = p_{ij}.
  \]
  $\mathbf{P} = (p_{ij})$: \textbf{ma trận chuyển}.
\end{dinhnghia}

\begin{tinhchat}
  \begin{itemize}
    \item $p_{ij} \ge 0$, $\sum_j p_{ij} = 1$ (hàng ngẫu nhiên).
    \item $P(X_n = j | X_0 = i) = (\mathbf{P}^n)_{ij}$ (xác suất chuyển $n$ bước).
    \item \textbf{Chapman--Kolmogorov:} $\mathbf{P}^{m+n} = \mathbf{P}^m \cdot \mathbf{P}^n$.
  \end{itemize}
\end{tinhchat}

\begin{dinhnghia}[title={Phân loại trạng thái}]
  \begin{itemize}
    \item $i \to j$ (thông nhau): $\exists n \ge 1$ sao cho $p_{ij}^{(n)} > 0$.
    \item \textbf{Bất khả quy}: mọi cặp trạng thái thông nhau.
    \item \textbf{Chu kỳ} $d(i) = \gcd\{n : p_{ii}^{(n)} > 0\}$. \textbf{Phi chu kỳ}: $d = 1$.
    \item \textbf{Hồi quy} (recurrent): $P(\text{quay lại } i | X_0 = i) = 1$.
    \item \textbf{Tạm thời} (transient): xác suất quay lại $< 1$.
    \item \textbf{Hồi quy dương}: $E[\text{thời gian quay lại}] < \infty$; \textbf{hồi quy null}: $= \infty$.
  \end{itemize}
\end{dinhnghia}

\begin{dinhly}[title={Phân phối dừng (stationary distribution)}]
  $\boldsymbol{\pi} = (\pi_j)$ là phân phối dừng nếu $\boldsymbol{\pi}\mathbf{P} = \boldsymbol{\pi}$, $\sum \pi_j = 1$, $\pi_j \ge 0$.

  Nếu xích bất khả quy, phi chu kỳ, hồi quy dương thì:
  \begin{enumerate}
    \item Tồn tại duy nhất $\boldsymbol{\pi}$.
    \item $\pi_j = 1/E_j[\tau_j]$ ($\tau_j$: thời gian quay lại $j$).
    \item $\lim_{n \to \infty} p_{ij}^{(n)} = \pi_j$ (hội tụ về phân phối dừng).
  \end{enumerate}
\end{dinhly}

\begin{dinhly}[title={Định lý Ergodic}]
  Dưới điều kiện trên: $\dfrac{1}{n}\sum_{k=0}^{n-1} f(X_k) \xrightarrow{\text{a.s.}} \sum_j f(j)\,\pi_j$.
\end{dinhly}

\begin{dinhnghia}[title={Cân bằng chi tiết (detailed balance)}]
  $\pi_i\, p_{ij} = \pi_j\, p_{ji}$ $\forall i, j$. Nếu thỏa thì $\boldsymbol{\pi}$ là phân phối dừng (xích khả đảo).
\end{dinhnghia}

%% ================================================================
%%  CHƯƠNG 3: QUÁ TRÌNH POISSON
%% ================================================================

\begin{tongket}[title={Tổng kết: Xích Markov}]
\begin{enumerate}
  \item Nắm lại các \textbf{định nghĩa} và \textbf{công thức} cốt lõi trong mục ``Xích Markov''.
  \item Tự làm lại ít nhất một ví dụ (nếu có) hoặc áp dụng công thức vào một tình huống số.
  \item Ghi chú điều kiện áp dụng (giả thiết độc lập, phân phối chuẩn, $n$ đủ lớn, v.v.).
\end{enumerate}
\end{tongket}

\section{Quá trình Poisson}

\begin{dongco}[title={Động cơ --- Quá trình Poisson}]
Đếm sự kiện rare trong thời gian liên tục.
\end{dongco}

\begin{dinhnghia}
  $\{N(t)\}_{t \ge 0}$ là \textbf{quá trình Poisson} cường độ $\lambda > 0$ nếu:
  \begin{enumerate}
    \item $N(0) = 0$.
    \item Gia số độc lập.
    \item $N(t) - N(s) \sim \text{Poisson}\bigl(\lambda(t-s)\bigr)$.
  \end{enumerate}
\end{dinhnghia}

\begin{tinhchat}
  \begin{itemize}
    \item $E[N(t)] = \lambda t$, $\text{Var}(N(t)) = \lambda t$.
    \item Thời gian giữa các sự kiện: $T_k \sim \text{Exp}(\lambda)$ iid.
    \item Thời điểm sự kiện thứ $n$: $S_n = \sum T_k \sim \text{Gamma}(n, \lambda)$.
    \item \textbf{Thinning}: lọc với xác suất $p$ $\Rightarrow$ Poisson$(\lambda p)$.
    \item \textbf{Chồng chất}: hợp hai Poisson độc lập $\lambda_1, \lambda_2$ $\Rightarrow$ Poisson$(\lambda_1 + \lambda_2)$.
  \end{itemize}
\end{tinhchat}

\begin{dinhnghia}[title={Quá trình Poisson không thuần nhất}]
  Cường độ $\lambda(t)$ thay đổi theo thời gian: $N(t) - N(s) \sim \text{Poisson}\!\left(\int_s^t \lambda(u)\,du\right)$.
\end{dinhnghia}

%% ================================================================
%%  CHƯƠNG 4: MARTINGALE
%% ================================================================

\begin{tongket}[title={Tổng kết: Quá trình Poisson}]
\begin{enumerate}
  \item Nắm lại các \textbf{định nghĩa} và \textbf{công thức} cốt lõi trong mục ``Quá trình Poisson''.
  \item Tự làm lại ít nhất một ví dụ (nếu có) hoặc áp dụng công thức vào một tình huống số.
  \item Ghi chú điều kiện áp dụng (giả thiết độc lập, phân phối chuẩn, $n$ đủ lớn, v.v.).
\end{enumerate}
\end{tongket}

\section{Martingale}

\begin{dongco}[title={Động cơ --- Martingale}]
``Trò chơi công bằng'': kỳ vọng tương lai bằng giá trị hiện tại.
\end{dongco}

\begin{dinhnghia}
  Quá trình $(M_n, \mathcal{F}_n)_{n \ge 0}$ là \textbf{martingale} nếu:
  \begin{enumerate}
    \item $M_n$ thích nghi với $\mathcal{F}_n$ (adapted).
    \item $E[|M_n|] < \infty$.
    \item $E[M_{n+1}|\mathcal{F}_n] = M_n$ a.s.
  \end{enumerate}
  Thay $=$ bằng $\le$: \textbf{supermartingale}. Thay bằng $\ge$: \textbf{submartingale}.
\end{dinhnghia}

\begin{vidu}
  \begin{itemize}
    \item Bước đi ngẫu nhiên đối xứng $S_n = \sum X_i$ ($P(X_i = \pm 1) = 1/2$).
    \item $M_n = S_n^2 - n$.
    \item Tỉ số hợp lý: $L_n = \prod_{i=1}^n \frac{f_1(X_i)}{f_0(X_i)}$.
    \item $M_n = e^{\theta S_n}/\bigl(E[e^{\theta X}]\bigr)^n$ (martingale mũ).
  \end{itemize}
\end{vidu}

\begin{dinhly}[title={Bất đẳng thức Doob}]
  Cho $(M_n)$ submartingale:
  \[
    P\!\left(\max_{0 \le k \le n} M_k \ge \lambda\right) \le \frac{E[M_n^+]}{\lambda}, \qquad \lambda > 0.
  \]
\end{dinhly}

\begin{dinhly}[title={Định lý hội tụ martingale (Doob)}]
  Nếu $(M_n)$ là martingale (hoặc submartingale) và $\sup_n E[M_n^+] < \infty$ thì $M_n \xrightarrow{\text{a.s.}} M_\infty$ hữu hạn.
\end{dinhly}

\begin{dinhly}[title={Optional Stopping Theorem}]
  Nếu $(M_n)$ martingale và $\tau$ là thời điểm dừng bị chặn ($\tau \le N$) thì $E[M_\tau] = E[M_0]$.

  Dạng mạnh hơn yêu cầu tính khả tích đều hoặc $E[\tau] < \infty$ kèm điều kiện bổ sung.
\end{dinhly}

%% ================================================================
%%  CHƯƠNG 5: CHUYỂN ĐỘNG BROWN
%% ================================================================

\begin{tongket}[title={Tổng kết: Martingale}]
\begin{enumerate}
  \item Nắm lại các \textbf{định nghĩa} và \textbf{công thức} cốt lõi trong mục ``Martingale''.
  \item Tự làm lại ít nhất một ví dụ (nếu có) hoặc áp dụng công thức vào một tình huống số.
  \item Ghi chú điều kiện áp dụng (giả thiết độc lập, phân phối chuẩn, $n$ đủ lớn, v.v.).
\end{enumerate}
\end{tongket}

\section{Chuyển động Brown (Wiener process)}

\begin{dongco}[title={Động cơ --- Chuyển động Brown (Wiener process)}]
Giới hạn của bước ngẫu nhiên nhỏ --- nền của tài chính và phương trình vi phân ngẫu nhiên.
\end{dongco}

\begin{dinhnghia}
  $\{B_t\}_{t \ge 0}$ là \textbf{chuyển động Brown chuẩn} nếu:
  \begin{enumerate}
    \item $B_0 = 0$.
    \item Gia số độc lập.
    \item $B_t - B_s \sim \mathcal{N}(0, t - s)$ ($0 \le s < t$).
    \item Quỹ đạo $t \mapsto B_t(\omega)$ liên tục a.s.
  \end{enumerate}
\end{dinhnghia}

\begin{tinhchat}
  \begin{itemize}
    \item $E[B_t] = 0$, $\text{Cov}(B_s, B_t) = \min(s, t)$.
    \item Quỹ đạo liên tục nhưng \textbf{không khả vi} tại bất kỳ điểm nào a.s.
    \item \textbf{Tự tương tự}: $\{c^{-1/2}B_{ct}\} \overset{d}{=} \{B_t\}$.
    \item \textbf{Tính Markov mạnh}: $\{B_{t+\tau} - B_\tau\}$ là Brown độc lập với $\mathcal{F}_\tau$.
    \item $B_t$ là martingale; $B_t^2 - t$ cũng là martingale.
  \end{itemize}
\end{tinhchat}

\begin{dinhly}[title={Biến phân bậc hai}]
  $[B]_t := \lim_{n \to \infty} \sum_{k} (B_{t_{k+1}} - B_{t_k})^2 = t$ a.s. (``$dB_t \cdot dB_t = dt$'').
\end{dinhly}

\begin{dinhnghia}[title={Chuyển động Brown hình học}]
  $S_t = S_0 \exp\!\left((\mu - \sigma^2/2)t + \sigma B_t\right)$, nghiệm SDE: $dS_t = \mu S_t\,dt + \sigma S_t\,dB_t$.
\end{dinhnghia}

\begin{luuy}
  Chuyển động Brown là nền tảng cho giải tích ngẫu nhiên (tích phân Itô) và tài chính toán (Black--Scholes).
\end{luuy}

\begin{tongket}[title={Tổng kết: Chuyển động Brown (Wiener process)}]
\begin{enumerate}
  \item Nắm lại các \textbf{định nghĩa} và \textbf{công thức} cốt lõi trong mục ``Chuyển động Brown (Wiener process)''.
  \item Tự làm lại ít nhất một ví dụ (nếu có) hoặc áp dụng công thức vào một tình huống số.
  \item Ghi chú điều kiện áp dụng (giả thiết độc lập, phân phối chuẩn, $n$ đủ lớn, v.v.).
\end{enumerate}
\end{tongket}



\end{document}
